% ======================================================================
\section{PLANET}
\label{sec:method}

\begin{figure*}[!t]
\centering
\figv{fig_arch_v2}{\textwidth}
\caption{PLANET. \textbf{(a)} The design matrix is a function of absolute local wall-clock
time, so it continues across the forecast origin without interruption: the phase of every
harmonic and the position of every weekend and holiday at horizon step $h$ are known exactly
before the forecast is made. \textbf{(b)} The phase-locked tier is solved in closed form per
window; a residual network models what it leaves behind, conditioned on the solved
coefficients; the shared exogenous head adds the covariate response. Gradients flow through
the solve.}
\label{fig:arch}
\end{figure*}

PLANET forecasts the load as the sum of a \emph{phase-locked} seasonal component re-estimated
for every window in closed form, a learned \emph{residual}, and a covariate \emph{response}
(Fig.~\ref{fig:arch}). The first is the contribution; the second is a backbone chosen on
validation; the third is the head every model in the study shares.

\subsection{Absolute-time bases}
\label{sec:basis}
For a local civil timestamp $\tau^{\ell}$ (hours) define the design vector
$\vv{\phi}(\tau^{\ell})\in\R^{K}$ by stacking an intercept; daily harmonics
$\{\sin,\cos\}(2\pi k\tau^{\ell}/24)$, $k=1{:}K_d$; weekly harmonics
$\{\sin,\cos\}(2\pi k\tau^{\ell}/168)$, $k=1{:}K_w$; and day-type terms---Saturday, Sunday and
public-holiday indicators, each alone and multiplied by the first two daily harmonic
pairs. The day-type group is what makes this a load-forecasting basis rather than a generic
Fourier basis: weekend and holiday demand is not weekday demand scaled down, it has a
different shape, and the interactions let the solver represent that at no extra cost. Two
optional components complete the seasonal representation: a local trend column of
$\vv{\phi}$, and a climatological tier $\vv{\Psi}\vv{\beta}^{g}_c$ of annual harmonics (with
an intercept and a multi-year trend) whose coefficients are learned globally and which is
removed from each channel before normalisation. The ablation of Section~\ref{sec:ablation}
measures what each contributes.

Two properties of $\vv{\phi}$ carry the method. It is a function of \emph{absolute} time: two
windows that end at the same wall-clock hour receive identical basis values wherever they
began, and a horizon step three weeks out is evaluated at its true calendar position, not at
an extrapolated offset. And it is evaluated on the \emph{local} clock: human schedules follow
civil time, so a basis on UTC is out of phase by one hour for half of every year in every
system studied here. Section~\ref{sec:ablation} measures both choices against their
alternatives---a window-relative basis, and a UTC basis---with everything else held fixed.

\subsection{The phase-locked least-squares solver}
Let $\vv{\Phi}_L\in\R^{L\times K}$ and $\vv{\Phi}_H\in\R^{H\times K}$ stack $\vv{\phi}$ over
the lookback and the horizon. The lookback, after removal of the climatological tier, is
instance-normalised per window and channel, $\hat{\mm{X}}=(\mm{X}-\vv{\mu})/\vv{\sigma}$,
which absorbs level and scale drift
\cite{kim2022revin}. For each window and channel $c$ the seasonal state is the solution of a
weighted ridge problem,
\begin{equation}
\begin{split}
\vv{\beta}_c&=\argmin_{\vv{b}}\;\sum_{i=1}^{L} w_{ic}\bigl(\hat{x}_{ic}-\vv{\phi}_i^\top\vv{b}\bigr)^2
+\sum_{k=1}^{K}\lambda_{ck}b_k^2\\
&=\bigl(\vv{\Phi}_L^\top\mm{W}_c\vv{\Phi}_L+\vv{\Lambda}_c\bigr)^{-1}\vv{\Phi}_L^\top\mm{W}_c\hat{\vv{x}}_c,
\end{split}
\label{eq:solution}
\end{equation}
with $\mm{W}_c=\operatorname{diag}(\vv{w}_{:c})$ and $\vv{\Lambda}_c=\operatorname{diag}(\vv{\lambda}_c)$,
and the seasonal forecast and its in-sample counterpart are
$\hat{\vv{s}}^{\text{fut}}_c=\vv{\Phi}_H\vv{\beta}_c$ and
$\hat{\vv{s}}^{\text{hist}}_c=\vv{\Phi}_L\vv{\beta}_c$.

Equation~\eqref{eq:solution} is a batch of symmetric positive-definite $K\times K$ systems with
$K\approx 35$, solved by Cholesky factorisation: there is no inner loop, no convergence
criterion and no possibility of divergence. The coefficients depend only on the current
window, so when the daily profile changes the very next forecast reflects it with the network
weights untouched---test-time adaptation obtained by solving a convex problem rather than by
continuing to train. And the solve is differentiable, so the weights $\vv{w}$, the penalties
$\vv{\lambda}$ and the residual network are trained for forecast error rather than for
in-sample fit.

The penalty is indexed by channel and basis function and parameterised multiplicatively,
$\lambda_{ck}=L\,(e^{\rho_{ck}}+10^{-3})$, with a larger step size on $\vv{\rho}$ than on the
network weights: its useful range spans several orders of magnitude, from almost none on a
clean daily profile to enough to erase the seasonal component of an aperiodic channel
entirely. The observation weights $\vv{w}_{:c}$ are either uniform or produced by a small
convolutional network with a learnable recency prior, normalised to $\sum_i w_{ic}=L$; the
headline model uses the learned weights and the ablation reports the uniform alternative.

\subsection{Residual network}
\label{sec:backbone}
What the tier leaves behind, $\mm{R}=\hat{\mm{X}}-\hat{\mm{S}}^{\text{hist}}$, is modelled by a
channel-independent network conditioned on the solved coefficients $\vv{\beta}_c$, which tell
it where in the calendar it is correcting and how strongly each harmonic is currently
expressed. Two backbones are candidates, and the choice between them is made once, on the validation loss
of the reference systems, and transferred like every other tuned setting
(Section~\ref{sec:setup}): the patch Transformer of PatchTST \cite{nie2023patchtst} (patches of 16 with stride 8, three encoder
layers, a flatten-and-project head), where $\vv{\beta}_c$ enters as one prepended token; and
the doubly-residual basis stack of N-BEATS \cite{oreshkin2020nbeats} (nine blocks of four
layers, width 256), where $\vv{\beta}_c$ is projected onto the input sequence and added to it
before the first block; a feature-wise affine modulation of every block's hidden state is the
alternative tested in the ablation. Both are the corresponding baseline architecture at the width selected for PLANET
(\appref{app:config}; the basis stack inside PLANET carries dropout, which the N-BEATS
baseline does not, and plain instance normalisation), and the ablation runs ``backbone
alone'' and ``backbone with the tier'' at that same width, so the pair is a controlled
measurement of what the tier contributes on two architectures that
share nothing else.

The point forecast is reassembled by inverting the normalisation,
$\hat{\vv{y}}=(\hat{\vv{s}}^{\text{fut}}_{c^\star}+\vv{r}^{\text{fut}}_{c^\star})\,\sigma_{c^\star}
+\mu_{c^\star}$. For probabilistic forecasts the residual head emits $Q$ quantile levels as a
median plus non-negative increments, trained with the pinball loss and calibrated as in
Section~\ref{sec:prob}.

\subsection{A native covariate pathway: PLANET-X}
\label{sec:planetx}
The shared head of Section~\ref{sec:head} is deliberately small, so that it corrects a
forecast rather than carries it, and PLANET receives the covariates through it like every
other model without a native interface. On systems whose load is driven by weather rather
than by the calendar, that pathway is a bottleneck that models with a native covariate
encoder do not have. PLANET-X is the same model with the head replaced by a native pathway of
the shape TiDE and N-BEATSx use for their exogenous inputs: a per-step linear map of the
covariate block over lookback and horizon to four features, flattened and projected to a
128-dimensional embedding, which is concatenated to the input of every block of the basis
stack (or prepended as one token to the patch Transformer). The calendar tier, the
normalisation and the conditioning on the seasonal state are unchanged. The variant was
designed after the first two systems' sweeps showed where the shared-head model loses, and we
say so; it is tuned and selected under the identical protocol and reported beside PLANET in
every experiment, and PLANET remains the primary model and the subject of the ablation.

\subsection{The shared exogenous head}
\label{sec:head}
Channel independence is what makes patch Transformers and basis stacks strong, but it also
means the covariate block of Section~\ref{sec:regimes} cannot influence the target through
them. Under our protocol every model without a native known-future interface receives the
block through one head, shaped like the classical load--temperature response: a point-wise
non-linearity applied to the covariates at every horizon step, a summary (mean, last, extrema)
of their lookback for thermal inertia, a short depth-wise temporal convolution, and a linear
output initialised near zero, so that a wrapped model starts exactly at its own forecast and
learns the covariate correction from there. The head is small
(\nv{proto.head.params_k}k parameters) and identical for every model it is attached to,
including PLANET.

\subsection{Complexity}
Building $\vv{\Phi}_L^\top\mm{W}\vv{\Phi}_L$ costs $O(BCLK^2)$ and the solves $O(BCK^3)$; with
$K\approx 35$ both are negligible beside either backbone. Under the load-only scoring of the
national grids every channel-independent model in the study, ours included, is trained on the
load channel alone, which yields the identical target forecast at a fraction of the cost;
multivariate models keep all channels. Measured cost is reported in
Section~\ref{sec:cost}.
